\documentclass[a4paper,12pt]{article}
\usepackage[T2A]{fontenc}
\usepackage[utf8]{inputenc}
\usepackage[russian]{babel}
\usepackage{amsmath,amssymb,mathtools}
\usepackage{helvet}
\usepackage{icomma}
\usepackage[a4paper,left=23mm,right=23mm,top=24mm,bottom=23mm,headheight=15pt]{geometry}
\usepackage{microtype}
\usepackage{tikz}
\usetikzlibrary{arrows.meta,positioning,shapes.geometric,calc}
\usepackage{pgfplots}
\pgfplotsset{compat=1.18}
\usepackage{tabularx,booktabs}
\usepackage{enumitem}
\usepackage[hidelinks]{hyperref}
\usepackage{parskip}
\usepackage{fancyhdr}
\usepackage{setspace}
\usepackage{xcolor}
\usepackage{titlesec}
\definecolor{accent}{HTML}{008080}
\definecolor{darkaccent}{HTML}{004D4D}
\definecolor{textgray}{HTML}{333333}
\definecolor{hair}{HTML}{A7C9C7}
\color{textgray}
\onehalfspacing
\setlength{\parindent}{0pt}
\setlist[itemize]{leftmargin=1.7em,itemsep=1pt,topsep=2pt}
\setlist[enumerate]{leftmargin=1.9em,itemsep=8pt,topsep=5pt}
\titleformat{\section}{\Large\color{darkaccent}}{\thesection.}{0.55em}{}[\vspace{1.5mm}\color{hair}\titlerule]
\titlespacing*{\section}{0pt}{1.05em}{0.6em}
\titleformat{\subsection}{\normalsize\bfseries\color{accent}}{}{0em}{}
\titlespacing*{\subsection}{0pt}{0.9em}{0.35em}
\pagestyle{fancy}
\fancyhf{}
\fancyhead[L]{\small Кинематика и движение тела}
\fancyhead[R]{\small Софья Лузянина}
\renewcommand{\headrulewidth}{0.35pt}
\fancyfoot[C]{\small\thepage}
\renewcommand{\footrulewidth}{0pt}
\newcommand{\keyidea}[1]{\noindent{\color{darkaccent}\textbf{Ключевая идея.}}~#1}
\newcommand{\vct}[1]{\boldsymbol{#1}}
\newcommand{\dd}{\mathop{}\!\mathrm{d}}
\begin{document}
\hypersetup{pageanchor=false}
\thispagestyle{empty}
\begin{center}
\vspace*{3.5cm}
{\Large\color{darkaccent} ФИЗИКА}\par
\vspace{1.45cm}
{\Huge\color{textgray} Чек-лист}\par
\vspace{0.75cm}
{\LARGE\color{accent} Кинематика: движение по осям,\\[2mm] горизонтальный и наклонный бросок}\par
\vspace{1.15cm}
{\normalsize Повторение методов и самостоятельная практика}\par
\vfill
{\large Софья Лузянина}\par
\vspace{0.7cm}
{16 ноября 2026 г.}\par
\vspace{1.4cm}
{\small Лёвин Артём Александрович}\par
{\small эксклюзивно для Софьи Лузяниной}\par
\vspace{1.15cm}
\begin{tikzpicture}[x=1cm,y=1cm]
\draw[accent,line width=.9pt] (-4.1,0) .. controls (-2.85,.85) and (-1.4,-.65) .. (0,0)
 .. controls (1.35,.72) and (2.8,-.54) .. (4.1,0);
\end{tikzpicture}
\end{center}
\clearpage
\hypersetup{pageanchor=true}
\setcounter{page}{1}

\section{Путь, перемещение и проекции}
\keyidea{Путь $l$ --- длина пройденной траектории (скаляр). Перемещение $\Delta\vct r$ --- вектор от начального положения к конечному. Всегда $l\ge |\Delta\vct r|$.}

Для движения по полуокружности радиуса $R$ между концами диаметра:
\[
 l=\pi R,\qquad |\Delta\vct r|=2R.
\]
После полного оборота путь равен $2\pi R$, а перемещение равно нулевому вектору.

\begin{center}
\begin{tikzpicture}[scale=.85,>=Stealth]
\draw[accent,line width=1.2pt,-{Stealth[length=2.4mm]}] (-2.25,0) arc[start angle=180,end angle=0,radius=2.25];
\draw[darkaccent,thick,-{Stealth[length=2.5mm]}] (-2.25,-.23)--(2.25,-.23);
\draw[gray,densely dashed] (0,0)--(0,2.25);
\fill (-2.25,0) circle (1.7pt) node[below left] {$A$};
\fill (2.25,0) circle (1.7pt) node[below right] {$B$};
\node[above right] at (.1,1.04) {$R$};
\node[below] at (0,-.35) {$\Delta\vct r$};
\node[above] at (-1.62,1.68) {траектория};
\end{tikzpicture}
\end{center}

Если вектор $\vct A$ образует угол $\alpha$ с положительной полуосью $Ox$, то его проекции:
\[
 A_x=A\cos\alpha,\qquad A_y=A\sin\alpha.
\]
Знаки проекций определяйте по направлениям осей: компонент в противоположном направлении получает знак минус.

\section{Уравнения движения и встреча тел}
При постоянном ускорении вдоль $Ox$:
\[
 x(t)=x_0+v_{0x}t+\frac{a_xt^2}{2},\qquad v_x(t)=v_{0x}+a_xt.
\]
Если $a_x=0$, движение равномерное: $x=x_0+v_xt$.

\subsection{Пример 1. Прочитать закон движения}
Для $x=5-4t+t^2$ (единицы СИ) получаем
\[
 x_0=5\ \text{м},\quad v_{0x}=-4\ \text{м/с},\quad a_x=2\ \text{м/с}^2,
\]
поскольку коэффициент при $t^2$ равен $a_x/2$.

\subsection{Пример 2. Найти встречу}
Пусть $x_1=5-4t$, $x_2=-3+2t$. В момент встречи координаты равны:
\[
 5-4t=-3+2t\ \Longrightarrow\ t=\frac43\ \text{с},\qquad
 x=5-4\cdot\frac43=-\frac13\ \text{м}.
\]
Время найдено относительно общего начала отсчёта; физический корень должен удовлетворять $t\ge0$.

\section{Графики и путь за отдельную секунду}
На графике зависимости $x(t)$ наклон прямой определяет \textit{проекцию} скорости:
\[
 v_x=\frac{x_2-x_1}{t_2-t_1}.
\]
Возрастание $x$ означает $v_x>0$; убывание --- $v_x<0$. При работе с графиком пользуйтесь числовыми шкалами осей.

\begin{center}
\begin{tikzpicture}
\begin{axis}[width=9.4cm,height=5.1cm,axis lines=middle,
 xlabel={$t,\ \text{с}$},ylabel={$x,\ \text{м}$},
 xmin=0,xmax=5.7,ymin=-5.5,ymax=12,
 xtick={0,1,2,3,4,5},ytick={-4,0,4,8,12},
 grid=major,grid style={draw=gray!16},
 unit vector ratio=1 1,clip=false,samples=201,
 every axis label/.append style={font=\small},tick label style={font=\small}]
\addplot[accent,very thick,domain=0:5] {10-3*x};
\addplot[only marks,mark=*,mark size=2.1pt,darkaccent] coordinates {(1,7) (5,-5)};
\end{axis}
\end{tikzpicture}
\end{center}

\textbf{За $n$-ю секунду.} Перемещение за интервал $[n-1;n]$ при $n\in\mathbb N$ и времени в секундах:
\[
 \Delta x_n=x(n)-x(n-1)=v_{0x}+\frac{a_x}{2}(2n-1).
\]
\textit{Путь} за этот интервал равен $|\Delta x_n|$ только при отсутствии смены направления. При развороте разделяйте интервал в момент $v_x=0$.

\subsection{Пример 3. Шестая секунда}
Тело движется в положительном направлении с $v_0=2$ м/с и $a=2$ м/с$^2$. Тогда
\[
 l_6=2+\frac{2}{2}(2\cdot6-1)=13\ \text{м}.
\]

\section{Бросок горизонтально с высоты}
Пусть начальная скорость $v_0>0$ направлена вдоль $+Ox$, ось $Oy$ направлена вверх, стартовая высота $h>0$; сопротивлением воздуха пренебрегаем, $g>0$.
\[
\begin{aligned}
 &a_x=0,\qquad v_x=v_0,\qquad &&x=v_0t,\\
 &a_y=-g,\qquad v_y=-gt,\qquad &&y=h-\frac{gt^2}{2}.
\end{aligned}
\]
В момент удара о горизонтальную площадку $y=0$:
\[
 0=h-\frac{gT^2}{2}\quad\Longrightarrow\quad
 \boxed{T=\sqrt{\frac{2h}{g}}}.
\]
Отсюда дальность и \textit{модуль} скорости удара:
\[
 \boxed{L=v_0\sqrt{\frac{2h}{g}}},\qquad
 \boxed{v_{\text{уд}}=\sqrt{v_0^2+2gh}}.
\]
Для текущего момента $t$:
\[
 v(t)=\sqrt{v_0^2+g^2t^2},\qquad
 \operatorname{tg}\varphi=\frac{gt}{v_0},
\]
где $\varphi\ge0$ --- угол \textit{ниже} горизонтали. Для ориентированного угла к $+Ox$ тангенс имеет знак минус.

\subsection{Пример 4. Время, дальность, скорость}
Мяч брошен горизонтально с $h=20$ м при $v_0=15$ м/с; $g=10$ м/с$^2$. Последовательно находим
\[
 T=\sqrt{40/10}=2\ \text{с},\qquad L=15\cdot2=30\ \text{м},
\]
\[
 v_{\text{уд}}=\sqrt{15^2+2\cdot10\cdot20}=25\ \text{м/с}.
\]
Важная проверка: $v_y(T)=-20$ м/с; $\sqrt{15^2+20^2}=25$ м/с.

\section{Бросок под углом к горизонту}
Пусть тело стартует с земли со скоростью $v_0$ под углом $0<\alpha<90^\circ$ к $+Ox$ и приземляется на \textit{тот же уровень}. Тогда
\[
 v_{0x}=v_0\cos\alpha,\qquad v_{0y}=v_0\sin\alpha;
\]
\[
 \begin{aligned}
 x(t)&=v_0\cos\alpha\,t, &v_x&=v_0\cos\alpha,\\
 y(t)&=v_0\sin\alpha\,t-\frac{gt^2}{2},
 &v_y&=v_0\sin\alpha-gt.
 \end{aligned}
\]
В вершине траектории $v_y=0$, значит
\[
 t_{\text{под}}=\frac{v_0\sin\alpha}{g},\qquad
 \boxed{H_{\max}=\frac{v_0^2\sin^2\alpha}{2g}}.
\]
Время полного полёта в указанном случае вдвое больше времени подъёма:
\[
 \boxed{T=\frac{2v_0\sin\alpha}{g}},\qquad
 \boxed{L=\frac{v_0^2\sin(2\alpha)}{g}}.
\]
Поскольку высота старта равна высоте приземления, $v_{\text{уд}}=v_0$ по модулю (без сопротивления воздуха).

\subsection{Пример 5. Четыре величины из двух данных}
При $v_0=20$ м/с, $\alpha=30^\circ$, $g=10$ м/с$^2$:
\[
 t_{\text{под}}=1\ \text{с},\quad T=2\ \text{с},\quad
 H_{\max}=5\ \text{м},\quad L=20\sqrt3\ \text{м}.
\]
Условия симметрии проверьте перед использованием формул\\
$T=2t_{\text{под}}$ и $L=v_0^2\sin(2\alpha)/g$.

\clearpage
\thispagestyle{fancy}
\begin{center}
{\LARGE\color{darkaccent} Алгоритм: движение тела в двух измерениях}\par
\vspace{1.4cm}
\end{center}
\begin{center}
\begin{tikzpicture}[>=Stealth,node distance=10mm and 17mm,
 action/.style={draw=accent,thick,rounded corners=2pt,align=center,minimum width=9.4cm,text width=9.05cm,minimum height=10mm,font=\small,inner sep=5pt},
 cond/.style={diamond,aspect=2.6,draw=darkaccent,thick,align=center,text width=4.6cm,inner sep=2pt,font=\small},
 side/.style={draw=accent,thick,rounded corners=2pt,align=center,text width=4.2cm,minimum height=13mm,font=\small,inner sep=4pt},
 term/.style={ellipse,draw=darkaccent,thick,align=center,minimum height=9mm,text width=4.6cm,font=\small},
 line/.style={-{Stealth[length=2.5mm]},thick,darkaccent}]
\node[term] (start) {Условие задачи};
\node[action,below=of start] (axes) {Задайте оси $Ox$, $Oy$ и их направления};
\node[action,below=of axes] (proj) {Разложите $\vct v_0$ на $v_{0x}$, $v_{0y}$; задайте $a_x$, $a_y$};
\node[action,below=of proj] (eq) {Запишите $x(t)$, $y(t)$, $v_x(t)$, $v_y(t)$ с правильными знаками};
\node[cond,below=14mm of eq] (choice) {Какое событие требуется найти?};
\node[side,below left=16mm and 10mm of choice] (apex) {Верхняя точка: $v_y=0$};
\node[side,below right=16mm and 10mm of choice] (impact) {Достижение уровня: $y=y_{\text{цель}}$};
\node[action,below=30mm of choice] (time) {Найдите время $t\ge0$ и проверьте его допустимость};
\node[action,below=of time] (calc) {Подставьте $t$ в $x(t)$, $y(t)$, $v_x(t)$ или $v_y(t)$};
\node[term,below=of calc] (end) {Проверьте ответ, знаки и единицы};
\draw[line] (start)--(axes);
\draw[line] (axes)--(proj);
\draw[line] (proj)--(eq);
\draw[line] (eq)--(choice);
\draw[line] (choice.south west) -| (apex.north);
\draw[line] (choice.south east) -| (impact.north);
\draw[line] (apex.south) -- ($(time.north west)!0.27!(time.north east)$);
\draw[line] (impact.south) -- ($(time.north west)!0.73!(time.north east)$);
\draw[line] (time)--(calc);
\draw[line] (calc)--(end);
\end{tikzpicture}
\end{center}
\clearpage

\section{Типичные ошибки и самопроверка}
\begin{itemize}
\item При чтении $x=x_0+v_{0x}t+\frac{a_xt^2}{2}$ не забывайте множитель $\frac12$: коэффициент при $t^2$ равен $a_x/2$.
\item Перемещение за промежуток времени может иметь знак; путь всегда неотрицателен и при развороте складывается по участкам.
\item Знак $v_y=-gt$ при горизонтальном броске следует из выбора оси $Oy$ вверх.
\item Для модуля скорости используйте теорему Пифагора: $v=\sqrt{v_x^2+v_y^2}$.
\item При ударе с высоты $h$ в выражении $v_{\text{уд}}^2=v_0^2+2gh$ стоит \textit{плюс}.
\item В верхней точке наклонного броска обнуляется $v_y$, тогда как $v_x$ остаётся постоянной.
\item Для броска под углом с приземлением на другой высоте время $T$ необходимо находить из $y(T)=y_{\text{цель}}$.
\end{itemize}

\subsection{Перед сдачей решения проверьте}
\begin{tabularx}{\linewidth}{@{}lX@{}}
$\square$ & Выбраны и подписаны направления осей.\\
$\square$ & Записаны проекции начальной скорости и ускорения.\\
$\square$ & Составлены уравнения движения в буквенном виде.\\
$\square$ & Указано условие искомого события: $y=0$ или $v_y=0$.\\
$\square$ & Числа подставлены после вывода нужного соотношения.\\
$\square$ & Сохранены размерности и выбран физический корень времени.\\
\end{tabularx}

\clearpage
\section*{Тренировочный блок --- путь А}
\addcontentsline{toc}{section}{Тренировочный блок}
\textit{10 заданий от базового до повышенного уровня. Для № 5--10 принимайте $g=10$ м/с$^2$; сопротивлением воздуха пренебрегайте. При решении задач о броске начните с анализа по осям.}

\subsection{Базовый уровень}
\begin{enumerate}[label=\textbf{\arabic*.}]
\item Тело прошло по полуокружности радиуса $R=3$ м из одного конца диаметра в другой. Найдите путь и модуль перемещения.
\item Координата тела задана формулой $x=4-6t+2t^2$ (СИ). Определите $x_0$, $v_{0x}$, $a_x$ и проекцию скорости при $t=2$ с.
\item Два тела движутся по одной прямой: $x_1=14-3t$, $x_2=-10+5t$ (СИ). Найдите момент и координату встречи.
\item Прямая на графике $x(t)$ проходит через точки $(t_1=1\ \text{с};\ x_1=7\ \text{м})$ и $(t_2=5\ \text{с};\ x_2=-5\ \text{м})$. Найдите $v_x$, $x_0$ и направление движения.
\end{enumerate}
\subsection{Средний уровень}
\begin{enumerate}[label=\textbf{\arabic*.},start=5]
\item Тело движется прямолинейно с $v_0=4$ м/с, $a=2$ м/с$^2$ в направлении скорости. Какой путь оно проходит \textit{за шестую секунду} движения?
\item Мяч брошен горизонтально с высоты $20$ м при $v_0=15$ м/с. Найдите время полёта, дальность и модуль скорости в момент падения.
\end{enumerate}
\clearpage
\subsection*{Тренировочный блок --- продолжение}
\begin{enumerate}[label=\textbf{\arabic*.},start=7]
\item С обрыва высотой $45$ м тело брошено горизонтально и упало на расстоянии $60$ м от основания. Найдите начальную скорость и тангенс угла, под которым скорость при ударе направлена ниже горизонта.
\item Мяч брошен горизонтально со скоростью $12$ м/с с высоты $80$ м. Найдите его координаты $x,y$ через $2$ с от броска (начало координат --- на земле под точкой броска), а также модуль скорости в этот момент.
\end{enumerate}
\subsection{Повышенный уровень}
\begin{enumerate}[label=\textbf{\arabic*.},start=9]
\item Тело брошено с земли под углом $45^\circ$ к горизонту со скоростью $20$ м/с и упало на тот же уровень. Найдите время подъёма, полное время полёта, максимальную высоту и дальность.
\item Тело брошено с земли под острым углом к горизонту и упало на тот же уровень. Через $2$ с после броска оно достигло вершины траектории, а горизонтальная проекция скорости в полёте была $15$ м/с. Найдите $v_0$, $\operatorname{tg}\alpha$, максимальную высоту и дальность полёта.
\end{enumerate}

\clearpage
\section*{Ответы к тренировочному блоку}
\noindent\begin{tabularx}{\linewidth}{@{}cX@{}}
\toprule
\textbf{№} & \textbf{Ответ}\\
\midrule
1 & $l=3\pi$ м; $|\Delta\vct r|=6$ м.\\
2 & $x_0=4$ м; $v_{0x}=-6$ м/с; $a_x=4$ м/с$^2$; $v_x(2)=2$ м/с.\\
3 & $t=3$ с; $x=5$ м.\\
4 & $v_x=-3$ м/с; $x_0=10$ м; против $+Ox$.\\
5 & $15$ м.\\
6 & $T=2$ с; $L=30$ м; $v_{\text{уд}}=25$ м/с.\\
7 & $T=3$ с; $v_0=20$ м/с; $\operatorname{tg}\varphi=\frac32$.\\
8 & $x=24$ м; $y=60$ м; $v=4\sqrt{34}$ м/с $\approx23{,}3$ м/с.\\
9 & $t_{\text{под}}=\sqrt2$ с; $T=2\sqrt2$ с; $H_{\max}=10$ м; $L=40$ м.\\
10 & $v_0=25$ м/с; $\operatorname{tg}\alpha=\frac43$; $H_{\max}=20$ м; $L=60$ м.\\
\bottomrule
\end{tabularx}

\vspace{1.3cm}
\noindent{\color{darkaccent}\textbf{Итоговая памятка}}\\[-1mm]
\begin{tabularx}{\linewidth}{@{}p{.37\linewidth}X@{}}
\toprule
\textbf{Что требуется} & \textbf{Ключ к решению}\\
\midrule
Встреча двух тел & Приравнять их координаты в один момент времени.\\
Путь за $n$-ю секунду & Вычесть координаты в $n$ и $n-1$; проверить разворот.\\
Момент удара о землю & Положить $y(T)=0$, выбрать $T>0$.\\
Вершина траектории & Положить $v_y=0$.\\
Дальность, высота, скорость & Найти нужное время и подставить в $x(t)$, $y(t)$, $v_x$, $v_y$.\\
\bottomrule
\end{tabularx}
\end{document}
